\section{Shared prefix conditioning}
\label{sec:prefix}

Continue to fix the feedback vertex set $F$ of cost at most $D$ and the
topological order $\prec$ from Section~\ref{sec:ranks}.  When we sample graph
tests or their defining data below, we use the finite laws that define the
graph.  In particular, the order is fixed before these
draws; the distribution of an embedding is not conditioned on whether any
test vertex survives.

We first specify a distribution that will be used at every bias scale.
Draw $\psi$ from the rank-spreading law.  Independently choose $t$ uniformly
from $\{0,\ldots,T-1\}$ and draw the prefix data of length $t$ from
Definition~\ref{def:prefix-data}.  Thus
$\mathbf s=(u_1,\ldots,u_t)$ has independent coordinates with law $\nu_U$,
the indices $i_1,\ldots,i_t$ are independent and uniform in $[M]$, and,
conditionally on these indices, the sets $P_j$ are independent with laws
$\mu_{p_{i_j}}^X$.  Write
\begin{equation}
  \mathcal D
  =\bigl(\psi,t,\mathbf s,(i_j,P_j)_{j=1}^t\bigr)
  \label{eq:common-conditioning}
\end{equation}
for these shared data.  Independently draw an edge $e=(u,v)$ with law
$\nu$.  Unless another experiment is explicitly specified, expectations
and probabilities involving sampled endpoints refer to this joint law of
$(\mathcal D,e)$.  In particular, its marginal law for $(\mathcal D,u)$
is obtained by drawing $u\sim\nu_U$ independently of $\mathcal D$; the
corresponding statement holds for $v\sim\nu_V$.

\begin{definition}[Conditioned functions]
\label{def:conditioned-functions}
For each game vertex $w$, fix a label
$\lambda_w^\circ\in\Lambda_w$ independently of all the sampled data.
Given $\mathcal D$, let $C_1,\ldots,C_{2^t}$ be its ordered prefix
partition.  If $\mathbf s$ is good, define its transition index by
\[
  r_*=\min\{r:g_{\mathbf s}(C_{\le r})=1\}.
\]
For every $w\in U\cup V$ define a function on $2^{\Lambda_w}$ by
\[
  f_w(R)=
  \begin{cases}
    g_{\mathbf s w}\bigl(E_w^{r_*}(R)\bigr),
      &\text{if $\mathbf s w$ is good},\\
    \ind{\lambda_w^\circ\in R},
      &\text{otherwise},
  \end{cases}
\]
where, as in Equation~\eqref{eq:comparison-sets},
\[
  E_w^r(R)
  =(C_{<r}\times\Lambda_w)\cup(C_r\times R).
\]
When the dependence on the shared data needs to be displayed, write
$f_{w,\mathcal D}$ in place of $f_w$.  For $0<p<1$, let
\[
  h_w(p)
  =\E_{R\sim\mu_p^{\Lambda_w}} f_w(R),
\]
and similarly write $h_{w,\mathcal D}(p)$ when needed.  This expectation
is taken with the function fixed.
\end{definition}

These definitions are valid also when some prefix cells are empty.
Indeed, on a good tuple the cumulative-set bits are monotone, start at
$0$, and end at $1$, by Lemma~\ref{lem:monotonicity}.  There is therefore
a unique transition, characterized by
\begin{equation}
  g_{\mathbf s}(C_{<r_*})=0,
  \qquad g_{\mathbf s}(C_{\le r_*})=1.
  \label{eq:prefix-transition}
\end{equation}
The transition cell is nonempty: adding an empty cell cannot change a
bit.  If $\mathbf s w$ is good, then $\mathbf s$ is good by
Lemma~\ref{lem:compatibility}, so the first case in
Definition~\ref{def:conditioned-functions} only uses a defined $r_*$.
If $\mathbf s$ is bad, all its extensions are bad and every function
uses the second case.

For a good extension, compatibility and
Equation~\eqref{eq:prefix-transition} give
\[
  f_w(\varnothing)=g_{\mathbf s}(C_{<r_*})=0,
  \qquad
  f_w(\Lambda_w)=g_{\mathbf s}(C_{\le r_*})=1.
\]
The set $E_w^{r_*}(R)$ increases with $R$, so this $f_w$ is monotone.
The fallback function has the same properties.  Consequently, for
\emph{every} realization of $\mathcal D$ and every vertex $w$, the
function $f_w$ is monotone with bottom $0$ and top $1$.  The entire
family $(f_w)_{w\in U\cup V}$ is determined by $\mathcal D$, without
using the fresh edge or a fresh scale index.

We will repeatedly use the marginal bounds
\begin{equation}
  \Pr[\mathbf s u\text{ is bad}]\le\rho,
  \qquad
  \Pr[\mathbf s v\text{ is bad}]\le\rho.
  \label{eq:sampled-extension-bad}
\end{equation}
To obtain them, condition on the sampled tuple and apply
Lemma~\ref{lem:goodness}; each tuple is independent of $\psi$.
This argument does not assert that all extensions are good
simultaneously.

Our objective is to prove $\E h_u(p_i)\ge p_i/4$ for every $i\in[M]$
under this same law.  To relate acceptance to the full-prefix tests,
we follow the transition cell through successive refinements: on a
good full tuple, its membership bit at each step will equal the
corresponding conditioned function value.  The deletion-cost bound
will make the counts of membership-one steps large with positive
probability; averaging over a uniformly chosen prefix length will
then give the desired acceptance bound.

\begin{lemma}[Refinement of the transition cell]
\label{lem:prefix-refinement}
Draw prefix data of length $T$ and an independent embedding $\psi$.
Write $\mathbf s_j=(u_1,\ldots,u_j)$, and let
$C^{(j)}_1,\ldots,C^{(j)}_{2^j}$ be the partition at depth $j$.
For $j=0$, this is the single cell $\Lambda_{\varnothing}$.
Suppose that $\mathbf s_T$ is good, and let $\tau_j$ be the
transition index at depth $j$.  Then the full transition cell is
nonempty and, for every $0\le j\le T$,
\[
  C^{(T)}_{\tau_T}
  \subseteq C^{(j)}_{\tau_j}\times X^{T-j}.
\]
Let $\mathcal D_{j-1}$ denote the shared data consisting of the same
embedding, the first $j-1$ prefix slots, and the deterministic length
$j-1$.  For each $1\le j\le T$, the membership bit of slot $j$
throughout the full transition cell is
\begin{equation}
  \ind{\lambda_j\in P_j}
  =f_{u_j,\mathcal D_{j-1}}(P_j)
  \qquad
  \text{for every }
  (\lambda_1,\ldots,\lambda_T)\in C^{(T)}_{\tau_T}.
  \label{eq:transition-membership}
\end{equation}
\end{lemma}

\begin{proof}
Goodness is inherited by all the prefixes.  A cumulative union of
cells at depth $j$, lifted to the full product alphabet, is a
cumulative union at a boundary of the depth-$T$ partition.  This
follows from the lexicographic ordering: all refinements of an
earlier cell precede all refinements of a later cell.  By
Lemma~\ref{lem:compatibility}, the cumulative-set bit at that lifted
boundary is its bit at depth $j$.  In particular, the two boundaries
on either side of $C^{(j)}_{\tau_j}$ have bits $0$ and $1$.
Monotonicity places the full transition between these boundaries,
which proves the asserted containment.  Its nonemptiness follows
because an empty cell cannot change a cumulative-set bit.

At step $j$, the depth-$(j-1)$ transition cell has two children:
its membership-$1$ child first, followed by its membership-$0$
child.  The cumulative union at the boundary between these children
is exactly
\[
  \bigl(C^{(j-1)}_{<\tau_{j-1}}\times X\bigr)
  \cup
  \bigl(C^{(j-1)}_{\tau_{j-1}}\times P_j\bigr).
\]
The bit of this set under $g_{\mathbf s_j}$ is
$f_{u_j,\mathcal D_{j-1}}(P_j)$.  Here the extension
$\mathbf s_{j-1}u_j=\mathbf s_j$ is good, so the fallback rule is
not used.  The bit before both children is $0$ and the bit after
both is $1$.  The transition is therefore in the first child if
and only if the intermediate bit is $1$.  If either child is
empty, the same statement follows from equality of its two
boundary sets.  The full transition refines the selected child,
giving Equation~\eqref{eq:transition-membership}.
\end{proof}

Figure~\ref{fig:prefix-refinement} illustrates the two-step case of the
refinement. Each displayed cell is a set of label tuples; its bit is
the value of $g$ on the cumulative union ending at that cell.
\begin{figure}[tb]
\centering
\begin{tikzpicture}[x=1.65cm,y=1cm,>=Stealth,
  cell/.style={draw,rounded corners=2pt,minimum width=1.35cm,
    minimum height=.65cm,font=\small},
  chosen/.style={cell,fill=blue!10,draw=blue!60!black}]
\node[cell] (root) at (0,2.6) {$\Lambda_{\varnothing}$};
\node[chosen] (one) at (-1.2,1.3) {$1$};
\node[cell] (zero) at (1.2,1.3) {$0$};
\node[cell] (oo) at (-1.8,0) {$11$};
\node[chosen] (oz) at (-.6,0) {$10$};
\node[cell] (zo) at (.6,0) {$01$};
\node[cell] (zz) at (1.8,0) {$00$};
\foreach \a/\b in {root/one,root/zero,one/oo,one/oz,zero/zo,zero/zz}
  \draw[->] (\a)--(\b);
\node[font=\small,anchor=west] at (2.4,1.3) {membership pattern};
\node[font=\small,anchor=west] at (2.4,-.75) {cumulative bit};
\foreach \x/\b in {-1.8/0,-.6/1,.6/1,1.8/1}
  \node[font=\small] at (\x,-.75) {$\b$};
\draw[->,blue!60!black] (-1.25,-.75)--(-1.05,-.75);
\end{tikzpicture}
\caption{Schematic transition refinement for two prefix slots. Membership
$1$ precedes $0$ lexicographically. In the illustrated case the cumulative
bit changes at cell $10$, so the transition is in the first cell at depth
one and its second child at depth two. The selected membership bits are
therefore $f_{u_1,\mathcal D_0}(P_1)=1$ and $f_{u_2,\mathcal D_1}(P_2)=0$, with each function using
its preceding prefix. Empty cells are retained in the ordering but cannot
contain a transition.}
\label{fig:prefix-refinement}
\end{figure}


The lemma uses formal product alphabets and literal prefix-lift
identities.  It does not require distinct game vertices in the
prefix.  In particular, a repeated vertex still occupies a
separate slot, with its own independently sampled index and
membership set.

\begin{proposition}[Average acceptance at every scale]
\label{prop:prefix-lower-bound}
Under the common law of $(\mathcal D,e)$, for every $i\in[M]$,
\begin{equation}
  \E h_u(p_i)\ge\frac{p_i}{4}.
  \label{eq:prefix-lower-bound}
\end{equation}
\end{proposition}

\begin{proof}
Consider the full-prefix experiment in
Definition~\ref{def:tests}, with data of length $T$ and an
independent embedding.  Denote its probability and expectation
by $\Pr_{\mathrm{full}}$ and $\E_{\mathrm{full}}$.
Let $Z_T$ be the star set of its wildcard test vertex, and let
$\mathcal A$ be the event that $\mathbf s_T$ is good and its
transition cell is not starred.

The star predicate is constant on each full cell: all label
tuples in that cell have the same membership vector, and hence
the same counts at each scale.  Suppose that $\mathbf s_T$ is
good and its transition cell is starred.  At the cut immediately
before that cell, put
\[
  E=C^{(T)}_{<\tau_T}\setminus Z_T.
\]
This is precisely the union of the nonstar portions of the cells
before the cut.  It is contained in $C^{(T)}_{<\tau_T}$, whose
bit is $0$.  Moreover, $E\cup Z_T$ contains
$C^{(T)}_{\le\tau_T}$, whose bit is $1$: adjoining the stars
restores every earlier starred cell and includes the entire
starred transition cell.  Monotonicity gives
\[
  g_{\mathbf s_T}(E)=0,
  \qquad g_{\mathbf s_T}(E\cup Z_T)=1.
\]
Lemma~\ref{lem:wildcard-forcing} forces deletion of this test
vertex.

The full-prefix family has total mass $H$, so its deletion
probability is at most $D/H$.  The full tuple is independent of
the embedding, and Lemma~\ref{lem:goodness} bounds its badness
probability by $\rho$.  We conclude that
\begin{equation}
  \Pr_{\mathrm{full}}(\mathcal A)
  \ge 1-\rho-\frac{D}{H}
  \ge\frac12.
  \label{eq:full-prefix-favorable}
\end{equation}
The last inequality follows from the parameter choices:
$\rho\le\gamma\le1/4000$ and $D/H<1/8$.

Define $\mathcal D_{j-1}$ from this experiment as in
Lemma~\ref{lem:prefix-refinement}.  For a fixed $i\in[M]$, set
\[
  Y_i=\sum_{j=1}^T
    \ind{i_j=i}\,
    f_{u_j,\mathcal D_{j-1}}(P_j).
\]
On $\mathcal A$, the transition cell is nonstarred, so the
definition of $Z_T$ and
Equation~\eqref{eq:transition-membership} give
$Y_i\ge p_iT/(2M)$.  On all other outcomes, $Y_i\ge0$,
including outcomes for which some conditioned functions use
the fallback rule.  It follows from
Equation~\eqref{eq:full-prefix-favorable} that
\begin{equation}
  \E_{\mathrm{full}}Y_i\ge\frac{p_iT}{4M}.
  \label{eq:full-prefix-count-lower}
\end{equation}

We now calculate this expectation without conditioning on
$\mathcal A$.  For each $j$, condition on the embedding, the
first $j-1$ prefix slots with all their data, and the new
vertex $u_j$.  The function
$f_{u_j,\mathcal D_{j-1}}$ is determined by these choices:
both the goodness test for $\mathbf s_j$ and any fallback
choice are independent of the fresh $i_j$ and $P_j$.
Since $i_j$ is uniform in $[M]$ and $P_j$ has law
$\mu_{p_i}^X$ conditional on $i_j=i$,
\[
  \E_{\mathrm{full}}\!\left[
    \ind{i_j=i} f_{u_j,\mathcal D_{j-1}}(P_j)
    \,\middle|\,
    \psi,(u_\ell,i_\ell,P_\ell)_{\ell<j},u_j
  \right]
  =\frac1M h_{u_j,\mathcal D_{j-1}}(p_i).
\]
The marginal law of the conditioning data and $u_j$ is the
common law of $(\mathcal D,u)$ conditional on $t=j-1$.
Indeed, the new vertex has law $\nu_U$ independently of
the earlier slots, just as does the big endpoint of the
independently sampled edge.  Therefore
\begin{align*}
  \E_{\mathrm{full}}Y_i
  &=\frac1M\sum_{j=1}^T
      \E\bigl[h_u(p_i)\mid t=j-1\bigr]\\
  &=\frac TM\E h_u(p_i).
\end{align*}
Together with Equation~\eqref{eq:full-prefix-count-lower},
this proves Equation~\eqref{eq:prefix-lower-bound}.
\end{proof}

The shared distribution in
Equation~\eqref{eq:common-conditioning} has not changed with
$i$.  The favorable full-chain event was used only to bound
a nonnegative random sum from below; it was not used to
condition the expectation that identifies this sum with
$\E h_u(p_i)$.
